\begin{derivation}[从\eqnrefs{eq:ew-onshell-mass-field-ct-dp11}到\eqnrefs{eq:sw-counterterm-massrelation-dp11}]
横向重整化逆传播子写成
\begin{equation*}
 [D_{V,T}(p^2)]^{-1}
 =p^2-m_V^2c^2+\Sigma_{VV,T}^{\rm loop}(p^2)
 -\delta m_V^2c^2+(p^2-m_V^2c^2)\delta Z_V.
\end{equation*}
极点仍位于 $p^2=m_V^2c^2$ 且留数为一，等价于
\begin{align*}
 0&=\Re\widehat\Sigma_{VV,T}(m_V^2c^2)
 \quad\Longrightarrow\quad
 \delta m_V^2c^2=\Re\Sigma_{VV,T}^{\rm loop}(m_V^2c^2),\\
 0&=\Re\widehat\Sigma'_{VV,T}(m_V^2c^2)
 \quad\Longrightarrow\quad
 \delta Z_V=-\Re\Sigma_{VV,T}^{{\rm loop}\,'}(m_V^2c^2),
\end{align*}
即正文\eqnrefs{eq:ew-onshell-mass-field-ct-dp11}。

在壳定义 $s_W^2=1-m_W^2/m_Z^2$ 的一阶变分直接给
\begin{align*}
 2s_W\delta s_W
 &= -\delta\!\left(\frac{m_W^2}{m_Z^2}\right)
 =-c_W^2\left(
 \frac{\delta m_W^2}{m_W^2}-\frac{\delta m_Z^2}{m_Z^2}
 \right),
\end{align*}
因此
\begin{equation*}
 \frac{\delta s_W}{s_W}
 =-\frac{c_W^2}{2s_W^2}
 \left(
 \frac{\delta m_W^2}{m_W^2}-\frac{\delta m_Z^2}{m_Z^2}
 \right),
\end{equation*}
得到正文\eqnrefs{eq:sw-counterterm-massrelation-dp11}。
\end{derivation}



\begin{derivation}[从\eqnrefs{eq:delta-rho-definition-dp11}到\eqnrefs{eq:delta-rho-doublet-dp11}]
记 $F_L=(f_1,f_2)_L^T$、$M_i=m_ic^2$。$\Delta\rho$ 只保留破坏守护 $SU(2)$ 的零动量自能差，因此先抽出共同的 $SU(2)_L$ 耦合并定义
\begin{equation*}
\mathcal I(a,b)
\equiv\mu^{4-d}\!\int\frac{\dd^d\ell_E}{(2\pi)^d}
\frac{\ell_E^2}{(\ell_E^2+a)(\ell_E^2+b)},
\qquad a=M_1^2,\quad b=M_2^2,
\end{equation*}
其中已经 Wick 旋转到欧氏能量动量，$\mu$ 是维数正规化尺度。这个标量积分来自左手流的 Dirac 迹。以异味带电流为例，Minkowski 分子为
\begin{align*}
T^{\mu\nu}_{12}
&=\Tr\!\left[\gamma^\mu P_L(\slashed\ell+M_2)
\gamma^\nu P_L(\slashed\ell+M_1)\right]\\
&=\frac12\Tr(\gamma^\mu\slashed\ell\gamma^\nu\slashed\ell)
=2\left(2\ell^\mu\ell^\nu-g^{\mu\nu}\ell^2\right).
\end{align*}
第二步用了 $P_L\gamma^\mu=\gamma^\mu P_R$：两个质量项之间含 $P_RP_L=0$，奇数个 $\gamma$ 的交叉项迹也为零。Lorentz 对称积分
\begin{equation*}
\ell^\mu\ell^\nu\longrightarrow\frac{g^{\mu\nu}}{d}\ell^2
\end{equation*}
把横向自能的零动量系数约化到 $\mathcal I$。暂时抽去两种自能共有的规范耦合、颜色多重度和外部归一化后，带电流的 $(1,2)$ 异味环与第三弱同位旋流的两个同味环组合成唯一的标量差
\begin{equation*}
\mathcal J(a,b)
\equiv
\mathcal I(a,b)-\frac12\mathcal I(a,a)-\frac12\mathcal I(b,b).
\end{equation*}
当 $a=b$ 时三项逐项给出 $\mathcal J(a,a)=0$；同时，大 $\ell_E$ 展开中的共同二次和对数紫外项在这个差中先行相消。

对标量积分作部分分式分解：

\begin{equation*}
\frac{\ell_E^2}{(\ell_E^2+a)(\ell_E^2+b)}
=-\frac{a}{b-a}\frac1{\ell_E^2+a}
+\frac{b}{b-a}\frac1{\ell_E^2+b}
\end{equation*}
并记
\begin{equation*}
A(a)\equiv\mu^{4-d}\int\frac{\dd^d\ell_E}{(2\pi)^d}\frac1{\ell_E^2+a}.
\end{equation*}
于是
\begin{equation*}
\mathcal I(a,b)=\frac{-aA(a)+bA(b)}{b-a},
\qquad
\mathcal I(a,a)=A(a)+aA'(a).
\end{equation*}
在 $d=4-2\epsilon$ 中，$A(a)$ 的展开只需要保留统一的极点/常数部分 $K_\epsilon$ 与质量对数：
\begin{equation*}
A(a)=\frac{a}{16\pi^2}
\left[K_\epsilon+1-\ln\frac{a}{\mu_E^2}\right]+O(\epsilon),
\end{equation*}
这里 $\mu_E$ 是与质量能量同量纲的重整化尺度。代入上述三项组合，$K_\epsilon$、$\ln\mu_E^2$ 和所有与 $a+b$ 成共同局域多项式的部分逐项消失，得到有限结果
\begin{align*}
\mathcal I(a,b)-\frac12\mathcal I(a,a)-\frac12\mathcal I(b,b)
&=\frac{1}{32\pi^2}
\frac{a^2-b^2-2ab\ln(a/b)}{a-b}\\
&=\frac{1}{32\pi^2}
\left[a+b-\frac{2ab}{a-b}\ln\frac ab\right].
\end{align*}

最后恢复弱耦合。带电流顶角给 $g/\sqrt2$，再除以 $m_W^2=g^2v_E^2/(4c^4)$；用能量真空标度 $v_E=v c^2$ 可把共同系数写成 $2/v_E^2$。因此
\begin{align*}
\Delta\rho_f
&=\frac{2N_c}{v_E^2}\mathcal J(a,b)\\
&=\frac{N_c}{16\pi^2v_E^2}
\left[a+b-\frac{2ab}{a-b}\ln\frac ab\right].
\end{align*}
由正文低能匹配 $G_F^{(E)}/\sqrt2=1/(2v_E^2)$，即
\begin{equation*}
\frac1{v_E^2}=\sqrt2\,G_F^{(E)},
\end{equation*}
于是
\begin{equation*}
\Delta\rho_f
=\frac{N_cG_F^{(E)}}{8\pi^2\sqrt2}
\left[
M_1^2+M_2^2-
\frac{2M_1^2M_2^2}{M_1^2-M_2^2}
\ln\frac{M_1^2}{M_2^2}
\right],
\end{equation*}
得到正文\eqnrefs{eq:delta-rho-doublet-dp11}。
\end{derivation}
